\cotefile{15}
\scanpage{120}\mspage{78}
and since $\Phi_k\ne 0$ in $Q'$, $f_nQ'=Q$, then $\frac{\Phi_{\unsure{0|O}}}{\Phi_k},\ldots,\frac{\Phi_T}{\Phi_k}$, and
with them $F_1$ and $F_2$ are regular in $Q'$; from the relations
$F_2(Q')=F_2(Q)\ne 0$ follows the regularity of $R$ in $Q'$.

Let us now show that, for large $n$, the mapping $f_n$
becomes one-to-one. The equation $\varphi^{\unsure{nn_e|\eta\eta_e|m_e}}=0$ will give a
cross-\unsure{'}section of the variety $V^{(n)}$ by some hyperplane, not
containing $V^{(n)}$. Evidently, this cross-section is the image
of $s-1$-dimensional subvariety $\mathscr E_{\varphi}$ of the variety $V$
(the equation $\mathscr E_{\varphi}$ in $U^{(1)}$ - the set $\varphi^{(i)}=0$).

According to the induction hypothesis, there exists
$m(\mathscr E_{\varphi})$ such that the elements of the group \handwritten{\unsure{$\checkmark$|of}}
$H^{\unsure{0|o}}(\mathscr E_{\varphi},O_{\mathscr E_{\varphi}}[mE']_{\mathscr E_{\varphi}})$ will give a biregular imbedding of
the variety $\mathscr E_{\varphi}$ in a projective space. From the exact
sequence (19) we obtain that this imbedding of the variety
$\mathscr E_{\varphi}$ is the restriction of the mapping $f_m\colon V\to V^{(m)}$ on
the subvariety $\mathscr E_{\varphi}$. We obtain, in particular that the
image $\mathscr E_{\varphi}$ by the mapping $f_m$ possesses a dimension $s-1$.
Hence it follows that the cross-section of the variety $V^{(m)}$
by a hyperplane, not containing $V^{(m)}$, possesses a dimension
$s-1$ and, consequently,
\[
\handwritten{\dim V^{(m)}}=s.
\]
% archive sheet number: 119
\scanpage{121}
poids et la structure de polarisation sur $G^\circ$
\uline{se déterminent mutuellement}.

\deleted{Donc} [ À voir : \unsure{l'un|l'ens.|le gpe} des autom de $(G^\circ_{\mathbb R},\varepsilon_{\mathbb R},j_{\mathbb R})$
voire $(G^\circ_{\deleted{\unsure{\mathbb R}}},\varepsilon,\mu)$, voire $(G,\varepsilon,\mu)$ opère-t-il
transitivement sur $\operatorname{Pol}(G)$ ? \deleted{\unsure{En d'autres|Et de même}} \inserted{Par exemple}
\deleted{\unsure{termes}, y a-t-il} \inserted{est-il \unsure{impossible|possible}} \deleted{\unsure{un sens}} \deleted{\unsure{l'on}} de distinguer \deleted{\unsure{avec} \illegible{1}} \inserted{une}
polarisation ou un paquet de polarisations
préférentielles ? ]

Prenons p. ex. une catégorie de Motifs
géométriques sur $\mathbb F_q$, $q=p^n$, \deleted{\unsure{et}} Alors \deleted{\unsure{dans}}
il y a une classe de conjugaison d'hom
\[
\deleted{\unsure{j}}i\colon G_{m\,\mathbb C}\longrightarrow G_{\mathbb C}
\]
distinguée, avec
\[
\bar i\,i=j,\qquad \varepsilon i=\mathrm{id}
\]
% slanted
\margin{N.B. $i$ et \unsure{ses} conjugués sur $\mathbb Q$, \deleted{\unsure{ceux des} \illegible{2}} \inserted{\deleted{\unsure{division}}}, si $[K:\mathbb Q]=n$ \unsure{engendrent} $\operatorname{Hom}(G_{m\,\mathbb C},G^\circ_{\mathbb C})$}
définie, \deleted{\unsure{par}} si $G$ commutatif : $M\subset L$ \inserted{$\underset{K\ \mathrm{tot.}\ \text{réel}}{\cup}$ \unsure{quadrique|quadratique} tot. imag.},
par la forme $\lambda\mapsto\frac1n v_{\mathfrak p}(\lambda)$, où $\mathfrak p$ est un
idéal premier de $L$ sur $p$. $[\,v_{\mathfrak p}(\lambda)v_{\mathfrak p}(\bar\lambda)=
v_{\mathfrak p}(\lambda\bar\lambda)=v_{\mathfrak p}(q^{j^*(\lambda)})=j^*(\lambda);\,v_{\mathfrak p}(q)=1\,]$.
% slanted
\margin{\deleted{\unsure{est|cet}} \uuline{\unsure{canulé|annulé|Calculer}} $\longrightarrow$ cette forme \deleted{\unsure{originale}} n'est pas tjrs à valeurs entières, même sur le \unsure{s-groupe|ss-groupe} de $M$ où $j=0$, ex : prendre $q=p\ne2$ $L=\mathbb Q(\sqrt p)$, tel que nb de classes de $L$ soit $1$, p. ex $p:3,7,11,\ldots$}
En fait, il convient de ``normaliser'' $G_{\mathbb C}$ choix
d'un tel $i$, (NB. Un autom de $G^\circ$ \inserted{\fbox{sur $\mathbb Q$!}} qui \deleted{induit}
respecte $i$ est l'identité !). La structure d'effectivité
sur $M$ est déterminée grâce à $i$, celle
des $M_*$ (\unsure{le cat. de $h$|la cat. de $h$|$K$ ext. de $\mathbb Q$}) en particulier de $M_{\mathbb R},M_{\mathbb C}$,
par la connaissance du système de ses conjugués
% archive sheet number: 120
\scanpage{122}\mspage{77}
or, what comes to the same thing, the values
\[
\handwritten{\frac{\unsure{\Phi|\phi}_{\unsure{0|o}}}{\Phi_k}(Q),\ldots\,.,\frac{\Phi_T}{\Phi_k}(Q)}
\]
from which it follows that
\[
\handwritten{\bigl(\Phi_0^{(i)}(Q'),\ldots\,.,\Phi_T^{(i)}(Q')\bigr)}=\handwritten{\bigl(x_0(Q),\ldots\,.,x_T(Q)\bigr)}
\]
\margin{p. 210}and
\[
\handwritten{Q}=\handwritten{f_nQ'}
\]

It is easy to show that $f_n$ is a continuous mapping
(in Zariski's topology) and even a regular one. The latter
follows from the fact that, in $Q\in V^{(n)}$ ($x_k\ne 0$ at $Q$),
any regular function $R$ has the form
\[
\handwritten{R}=\handwritten{\frac{F_1\left(\frac{x_0}{x_k},\ldots,\frac{x_T}{x_k}\right)}{F_2\left(\frac{x_0}{x_k},\ldots\,.\frac{x_T}{x_k}\right)}}
\]
where $F_1$ and $F_2$ are polynomials and $F_2(Q)\ne0$; the
corresponding function on $V$ will be
\[
\handwritten{R'}=\handwritten{\frac{F_1\left(\frac{\Phi_0}{\Phi_k},\ldots\,.,\frac{\Phi_T}{\Phi_k}\right)}{F_2\left(\frac{\Phi_0}{\Phi_k},\ldots\,.,\frac{\Phi_T}{\Phi_k}\right)}}
\]
% archive sheet number: 121
\scanpage{123}
de $i$. \deleted{Question : la structure de positivité est elle \unsure{associée|analogie} à $i|_{\mu_2}$ ?} Donc, grâce à cette circonstance
spéciale, il y a des choix de \deleted{\unsure{polarisati|positivité}} \unsure{structure|structures}
polarisantes privilégiés \unsure{de|dans|du} $G$ : les
$i|_{\mu_2}$ avec $i$ dans la classe ci-dessus.
\margin{\diagram{\begin{tikzcd}[ampersand replacement=\&]G_{m\,\mathbb C}\arrow[r,dash,"i"]\& G_{\mathbb C}\\G_{m\,\mathbb C}\arrow[r,"i'"]\& G_{\mathbb C}\end{tikzcd}}}
Quel est l'arbitraire ? On constate que l'ens.
\inserted{des} $i-i'$ \deleted{\unsure{correspond|conjugués} : $v_{\mathfrak p}(\lambda)\,v_{\mathfrak p'}(\lambda)$} définit un
sous-groupe \uline{d'indice fini} de $\operatorname{Hom}(G_{m\,\mathbb C},G^\circ_{\mathbb C})$,
mais cela ne donne aucune indication sur la
\unsure{composante|comportement} des $i_0-i'_0$, où l'indice $0$ désigne le
\unsure{restr.|contr.} à $\mu_2$. Question semble secondaire.

\margin{$\vert\!\vert\!\vert$}En tous cas, on a \deleted{\unsure{des|deux}} \inserted{un} paquet de
$\mathbb R$-foncteurs fibres privilégiés, définis en termes des
différents $i$.

\begin{center}\rule{2.2cm}{0.3pt}\end{center}

Autre structure spécifique sur le $G$ \unsure{provenant|pour}
d'un corps fini : d'après la façon de
décrire $G_{\mathbb C}=D_{\mathbb C}(M)$, $M\subset L^*$, comme $L$ est
stable par conjugaison complexe, qui commute
à $\operatorname{Gal}(L/\mathbb Q)$, on trouve dans $M$ donc
$G$ une involution canonique, notée
\[
u\longmapsto u'
\]
Dans \deleted{\unsure{D}}$\operatorname{Ker}j^*$ c'est la symétrie, d'autre part sur $\operatorname{Im}j$ c'est
l'identité, ce qui suffit pratiquement à la déterminer. Sur
$\operatorname{Hom}(G_{m,\mathbb C},G_{\mathbb C})$ $u\mapsto u'$ induit la conjugaison complexe.
% archive sheet number: 122
\scanpage{124}\mspage{76}
of the field $k(\frac{x_{\unsure{0|o}}}{x_k},\ldots,\frac{x_T}{x_k})$ in $(k,\infty)$ and even the
homomorphism of the field $k(\frac{\Phi_0}{\Phi_k},\ldots,\frac{\Phi_T}{\Phi_k})$ in $(k,\infty)$. The
last homomorphism can be extended to the point of the
field $k(V)$, and, in view of the completeness of $V$,
this\unsure{,} point possesses some center $Q'$ on $V$. Let $k_1$
be such that $\Phi_{k_1}\ne0$ in $Q'$. Then $F=\frac{\Phi_k}{\Phi_{k_1}}$ is a
regular function in $Q'$
(since, if $Q'\in U^{(i)}$, then $F=\frac{\Phi_k^{(i)}}{\Phi_{k_1}^{(i)}}$ and $\Phi_k^{(i)}$,
$\Phi_{k_1}^{(i)}\in O_{V,Q'}$, $\Phi_{k_1}^{(i)}(Q')\ne0$). If $\Phi_k=0$ in $Q'$,
then $F(Q')=0$ and
\[
\handwritten{\frac{1}{F(Q')}}=\handwritten{\frac{\unsure{\Phi|\phi}_{k_1}}{\Phi_k}(Q')}=\infty
\]
But
\[
\handwritten{\frac{\Phi_{k_1}}{\Phi_k}(Q')=\frac{\Phi_{k_1}}{\Phi_k}(Q)}=\handwritten{\frac{x_{k_1}}{x_k}(Q)}\ne\infty
\]

Therefore, $\Phi_k\ne0$ in $Q'$. But then the functions
\[
\handwritten{\frac{\Phi_0}{\Phi_k},\ldots\mathord{.},\frac{\Phi_T}{\Phi_k}}
\]
are regular in $Q'$ and take there the values
\[
\handwritten{\frac{x_0}{x_k}(Q),\ldots\mathord{.},\frac{x_T}{x_k}(Q)}
\]
% archive sheet number: 123
\scanpage{125}
Il ne semble pas que cette involution corresponde
à une involution \unsure{dans|de|sur} la \uline{gerbe} $\unsure{\mathscr G|\mathcal G}$ associée à
$M$ (on \deleted{\unsure{n'aurait} \unsure{pas|tenant}} n'a pas la propriété de
covariance voulue p.r. à l'élément fondamental
correspondant \unsure{de|dans|des} $H^{\unsure{2|n}}(\mathbb Q,\deleted{\unsure{G}}U)$ !, qui
serait \deleted{\unsure{d'ordre|d'autres} \unsure{de}} d'ordre divisant $2\ldots$).
% archive sheet number: 124
\scanpage{126}\mspage{74}
% paper clip at upper right
\[
\handwritten{H^1(\mathscr E_{\varphi},O_{\mathscr E_{\varphi}}[nE]_{\mathscr E_{\varphi}})}=\handwritten{0}
\]
we obtain an exact sequence
\[
\handwritten{H^{\unsure{0|o}}(V,O_V[nE])\xrightarrow{r^*}H^0(\mathscr E_{\varphi},O_{\mathscr E_{\varphi}}[nE]_{\mathscr E_{\varphi}})\to0}\tag*{\handwritten{(19)}}
\]

Let us assume that there would exist a point $P\in\mathscr E_{\varphi}$
such that in $P$ any cross-section $\Phi\in H^0(V,O_V[nE])$
satisfies the condition: if $P\in U^{(i)}$, then $\Phi^{(i)}(P)=0$.
Since any cross-section $\Phi'\in H^0(\handwritten{\mathscr E_{\varphi}},\handwritten{O}_{\deleted{\unsure{\varphi|o}}\unsure{\mathscr E_{\varphi}|\varphi}}[nE]_{\mathscr E_{\varphi}})$ can be
presented in the form $\Phi'=r^*\Phi$ or $\Phi^{\prime(i)}=r\,\Phi^{(i)}$ and
$r$ in the stalk over $P$ is a homomorphism of restriction
of functions; then $\Phi^{\prime(i)}(P)=0$. But this contradicts
the fact that, at $n\ge n(\mathscr E_{\varphi})$, the elements of the group
$H^0(\mathscr E_{\varphi},O_{\mathscr E_{\varphi}}[nE]_{\mathscr E_{\varphi}})$ generate the stalk of the sheaf
$O_{\mathscr E_{\varphi}}[nE]_{\mathscr E_{\varphi}}$ over any point $P\in\mathscr E_{\varphi}$. The lemma is proved. $\blacksquare$

Let us come back to the proof of theorem 2. Put
$E'=n_{\mathscr E}E$, where $n_{\mathscr E}$ is the number to which lemma 3
refers. Let $\Phi_{\unsure{0|o}},\ldots,\Phi_T$ be the basis (not necessarily
of linearly independent vectors) of the group
$H^0(V,O_V[nE'])$.

Consider a mapping $f_n\colon V\to P^T$ ($P^T$ is a
$T$-dimensional projective space), operating according to
the rule: if $P\in V$, $P\in U^{(i)}$, then
% archive sheet number: 125
\scanpage{128}
\uline{Définition} Soit $K/\mathbb Q$ une extension \deleted{\unsure{alg. de gr.|de degré fini}} \inserted{algébrique}. On dit
que $K$ est tot. réel si pour tout $\mathbb Q$-hom $\sigma:K\to\mathbb C$,
on a $\sigma(K)\subset\mathbb R$. On dit qu'un élément $\unsure{a|x}$ de $K$ est
tot. réel \inserted{(resp. tot. $>0$, tot. $\ge0$)} si pour tt $\sigma$ comme dessus, on a
$\sigma(a)\in\mathbb R$ (resp. $\sigma(x)>0$, resp. $\sigma(a)\ge0$). Tot.\textsuperscript{t} (ou purement)
imaginaire : \ldots .

\uline{Sorite} Pour que $x$ soit tot. réel, il f. et s. que
$\mathbb Q(x)\subset K$ le soit. Pour qu'il soit tot. $\ge0$, il faut et
il suffit qu'il soit nul ou tot $>0$. Les éléments
tot $\ge0$ \deleted{\unsure{forment}} \inserted{sont} les éléments $\ge0$ pour une relation
d'ordre sur $K$. [On parlera de $x$ tot. $\ge y$ ou
$x$ tot. $\le y$]. Si $K$ est engendré par des ss-corps
$K_\alpha$, alors $K$ est tot. réel ssi les $K_\alpha$ le sont. Si
\inserted{$K$ est engendré par des $x_\alpha$,}
$K$ tot. réel ssi \deleted{\unsure{tous les}} les $x_\alpha$ le sont.

Il existe, pour $K$ donné, un plus grand ss-corps
totalement réel de $K$, appelé la \uline{clôture tot. réelle} de
$\mathbb Q$ dans $K$. Soit $\deleted{\unsure{K=|K}}\overline{\mathbb Q}\subset\mathbb C$, et $\pi=\operatorname{Gal}(\deleted{\unsure{K}}\overline{\mathbb Q}/\mathbb Q)$,
enfin $\tau\in\pi$ la conjugaison complexe
\[
\tau\lambda=\bar\lambda\qquad \text{pour }\lambda\in\overline{\mathbb Q}
\]
de sorte qu'on a
\[
\tau^2=1.
\]

Alors le ss-groupe $\pi'$ de $\pi$ qui \deleted{\unsure{invarie|invariant}} \inserted{correspond} à clôture
tot\textsuperscript{t} réelle $\deleted{\unsure{de|K}}\mathbb K\ \deleted{\unsure{\mathbb Q}}$ de $\mathbb Q$ dans $\overline{\mathbb Q}$ est \deleted{\unsure{le}} le ss-groupe
\margin{\diagram{\begin{tikzcd}[ampersand replacement=\&,row sep=small]A\arrow[d,dash]\\\unsure{L|\mathbb L}\arrow[d,dash]\\\unsure{K|\mathbb K}\arrow[d,dash]\\\mathbb Q\end{tikzcd}}}
\uline{invariant} \uline{fermé} engendré par $\tau$, i.e. le sous groupe invariant
engendré par les $ {}^{\sigma}\tau=\sigma\tau\sigma^{-1}$ $(\sigma\in\pi)$.

Considérons la plus grande extension abélienne
$A$ de $\mathbb K$, dont le groupe de Galois est donc $H\unsure{=|:}\pi'/(\pi',\pi')$. \deleted{\unsure{On|Son}} Sur
% archive sheet number: 127
\scanpage{129}
ce groupe $G=\pi/\pi'=\operatorname{Gal}(\mathbb K/\mathbb Q)$ opère continûment, et
il est engendré topologiquement par les $ {}^\sigma\dot\tau$, où $\dot\tau\in H$
est l'image de $\tau$, et $\sigma\in G$. Donc il \uline{est annulé par} $2$.

\margin{$\left[\rule{0pt}{9em}\right.$}Soit $L$ une extension abélienne de $\mathbb K$, \deleted{\unsure{envisag|correspondant} \unsure{les conditions}} \inserted{dont le}
groupe de Galois est donc un quotient $\Gamma$ de $H$; \deleted{Conditions}
Supposons $L\ne\mathbb K$, i.e. $\Gamma\ne\{e\}$. Conditions équivalentes

(i) les $\sigma\tau\sigma^{-1}$ $(\sigma\in\pi)$ coïncident toutes sur $L$
(donc y coïncident avec $\tau:\lambda\longrightarrow\bar\lambda$)

(ii) $L$ est une extension \uline{quadratique} de $\mathbb K$ \uline{galoisienne}
\uline{sur} $\mathbb Q$.

\deleted{\unsure{L est une extension} \illegible{2}} \inserted{(iii) $L$ est une extension \uline{quadratique} \uline{purement imaginaire} de $\mathbb K$.}

De plus, une telle $L$ \deleted{existe et} est unique, [et
d'ailleurs existe : il suffit de prendre une extension
quadratique imaginaire $E$ de $\mathbb Q$\inserted{(p. ex. $\mathbb Q(\sqrt{-1})$)}, et $L=\mathbb Q(\mathbb K,E)$
(p. ex. $L=\mathbb K(\sqrt{-1})$)].

\uline{Dém} $\mathrm{(i)}\Rightarrow\mathrm{(ii)}$ Comme $\Gamma$ est engendré par les images
\deleted{$\unsure{{}^\sigma\dot\tau}$} des $ {}^\sigma\tau$, \deleted{\unsure{il s'}} et que l'hyp. de (i) signifie que
ces images sont égales, on voit que (i) implique que
$\Gamma$ est engendré par l'élément $\dot\tau$ de degré $\mid 2$, donc est
$\simeq\mathbb Z/2\mathbb Z$ (car $\Gamma\ne(e)$ par hypothèse). D'ailleurs, si un
tel $\Gamma$ existe, \deleted{\unsure{on}} \inserted{il est} nécessairement égal à $H_G$ (car \deleted{\unsure{l'on}}
c'est un quotient de $H_G$, et l'argument précédent montre
que \deleted{\unsure{ce}} $H_G$ est engendré par un élément d'ordre $\mid2$. \deleted{\unsure{On}}
\deleted{il existe bel et} Il en résulte aussitôt que $L$ est galoisien
sur $\mathbb Q$, donc $\mathrm{(i)}\Rightarrow\mathrm{(ii)}$. D'autre part, $\mathrm{(ii)}\Rightarrow\mathrm{(i)}$, car
(ii) implique que $\Gamma\simeq\mathbb Z/2\mathbb Z$\inserted{ et que $G$ y opère}, et comme $\unsure{\Gamma|F}$ est engendré par les
images des $\sigma\tau\sigma^{-1}={}^{\sigma}\tau$, i.e. par les $ {}^\sigma\dot\tau$, on voit que
$\dot\tau$ (donc les $ {}^\sigma\dot\tau$) sont tous égaux à l'unique élément non nul de $\Gamma$.
$(\mathrm{(i)},\mathrm{(ii)})\Longleftrightarrow\mathrm{(iii)}$ car du moment que $L$ est quadratique sur $\mathbb K$,
la condition (i) sur les $\sigma\tau\sigma^{-1}$ est \unsure{évidemment|évidemt} équival\deleted{\unsure{ente}} à (iii)
% archive sheet number: 128
\scanpage{130}
\uline{Corollaire} Soit $L$ un sous corps de $\overline{\mathbb Q}$. Pour qu'on ait
$L\subset\mathbb L$, il f. et s. que $L$ soit tot\textsuperscript{t} réel, ou extension
quadratique purement imaginaire d'un corps tot. réel.

Suffisance triviale par (iii). Nécessité : Par théorie
de Galois, $L$ est une extension de degré $1$ ou $2$
\margin{\diagram{\begin{tikzcd}[ampersand replacement=\&,row sep=small,column sep=small]\mathbb K\arrow[r,dash]\arrow[d,dash]\&\mathbb K(L)\arrow[r,dash]\arrow[d,dash]\&\mathbb L\\K=\mathbb K\cap\unsure{\mathbb L|L}\arrow[r,dash]\arrow[d,dash]\&L\&\\\mathbb Q\&\&\end{tikzcd}}}
du corps tot. réel $K=L\cap\mathbb K$, et il reste à
prouver que lorsque $L$ en est une extension
quadratique, cette dernière est tot\textsuperscript{t} imaginaire.
\deleted{\unsure{Or on}} Or \deleted{\unsure{ceci résulte du}} on est ramené à prouver
\deleted{\uline{Lemme} Soit $L$ une extension quadratique d'un corps tot.
réel $K$. Si $L$}
(quitte à remplacer $L$ par un conjugué, \inserted{ce qui ne dérange pas $\mathbb K$},) que $L$ est
imaginaire. Or s'il était réel, comme $\mathbb K$ l'est, $\mathbb K(L)$ le
serait, or $\mathbb K(L)=\mathbb L$, absurde.

\uline{NB} $\mathbb L$ est la plus gde extension galoisienne de $\mathbb Q$ telle
que dans son gpe de Galois $G'$, l'image de $\tau$ (conjugaison complexe) soit \uline{centrale}. Donc $\mathbb L$ contient l'extension \deleted{\unsure{g|gal}}\unsure{abélienne|galoisienne}
maximale de $\mathbb Q$, i.e. le corps des
racines de l'unité.
% archive sheet number: 129
\scanpage{131}
\uline{Proposition} Soit $Q\in\mathbb Q^*$, et soit $\lambda\in\mathbb C$ un nombre algébrique.
Conditions équivalentes

(a) Pour tout conjugué $\lambda_\alpha$ de $\lambda$, on a $\lambda_\alpha\bar\lambda_\alpha=Q$ i.e. $|\lambda_\alpha|^2=\unsure{Q}$

(b) \deleted{\unsure{Il est équivalent}} Il existe un corps \inserted{de nbres} \uline{totalement réel} $K$, tel que
\deleted{\unsure{son s}} $K(\lambda)$ soit égal à $K$ ou une extension \uline{quadratique} purement
imaginaire de $K$, et on a $\lambda\bar\lambda=Q$ i.e. $|\lambda|^2=Q$.

(b') Comme (b), avec $K=\mathbb Q(a)\deleted{\unsure{=K\cap\mathbb R}}$, où $a=\lambda+\bar\lambda=\lambda+\frac Q\lambda$.

(b'') Il existe un nombre algébrique \deleted{\unsure{et|$a$}} \uline{tot} \uline{réel} $a$,
avec \deleted{$\unsure{a^2-4Q}$ nul} $4Q-a^2$ nul ou totalement
positif, tel que $\lambda$ satisfasse une équation
\[
\lambda^2-a\lambda+Q=0
\]

(c) Soit $\mathbb K$ l'extension alg. tot. réelle maximale de $\mathbb Q$, et
$\mathbb L$ l'unique extension quadratique de $\mathbb K$ galoisienne sur $\mathbb Q$.
On a $\lambda\in\mathbb L$, et \deleted{$\unsure{|\lambda|}$} $\lambda\bar\lambda=Q$ i.e. $|\lambda|^2=Q$.

\uline{Dém}

a) $\Longrightarrow$ b'') Posons $a=\lambda+\bar\lambda=\lambda+\frac Q\lambda$. Je dis que $a$ \uline{est tot. réel}.
En effet, les conjugués de $a$ sont les $\deleted{\unsure{\sigma}}a_\alpha=\deleted{\unsure{\sigma}}\lambda_\alpha+\frac{Q}{\deleted{\unsure{\sigma}}\lambda_\alpha}$, or
par hypothèse $\lambda_\alpha\bar\lambda_\alpha=Q$ donc $a_\alpha=\lambda_\alpha+\bar\lambda_\alpha\in\mathbb R$. D'autre part,
$\lambda$ est solution de l'équation
\[
(T-\lambda)(T-\bar\lambda)=0\qquad\text{i.e.}\qquad T^2-aT+Q=0,
\]
et par suite $\lambda_\alpha$ est solution de l'équation $T^2-a_\alpha T+Q_{\deleted{\unsure{q}}}=0$,
\deleted{de plus} l'autre solution étant $\frac Q{\lambda_\alpha}$\inserted{(qui est donc égale à $\bar\lambda_\alpha$)}. \deleted{\unsure{Donc}} Comme
$\lambda_\alpha\bar\lambda_\alpha=Q$, les deux solutions sont $\lambda_\alpha$ et $\bar\lambda_\alpha$, donc complexes
conjuguées. Cela montre que $a_\alpha^2-4Q\le0$, \deleted{\unsure{donc|d'où}} \inserted{i.e.}
$4Q-a_\alpha^2\ge0$, donc $4Q-a^2$ est tot positif.

b'') $\Longrightarrow$ b') Il est évident que $\lambda$ est de degré au plus
deux sur \inserted{le corps tot. réel} $\mathbb Q(a)$, \deleted{De plus \unsure{comme} $\unsure{a=\lambda+\frac Q\lambda}$} \inserted{et $\mathbb Q(\lambda)\supset\mathbb Q(a)$ [car $a=\lambda+\frac Q\lambda$].}
\deleted{\unsure{ce qui}} Comme $a^2-4Q\le0$, les solutions de l'équation
$T^2-aT+Q=0$ sont $\lambda$ et $\bar\lambda$\inserted{(cf lemme plus bas)}, et on a $\lambda\bar\lambda=Q$, $\lambda\unsure{+|-}\bar\lambda=a$.
% archive sheet number: 130
\scanpage{132}
\deleted{Évident} Il faut prouver que si $\lambda\notin\mathbb K$ i.e. $a^2-4Q\ne0$ (cf lemme
plus bas), alors \deleted{il est pure} $\mathbb Q(\lambda)$ est purement imaginaire, i.e.
les $\lambda_\alpha$ sont tous non réels. \deleted{Cela signifie que les} \inserted{Comme $\lambda_\alpha$ est solution}
de l'équation $T^{\unsure{2}}-a_\alpha T+Q=0$, cela signifie que
$a_\alpha^2-4Q<0$, ce qui résulte du fait que $4Q-a^2$ est tot $\ge0$.

b') $\Rightarrow$ b) trivial, \deleted{\unsure{puisque}}

b) $\Longrightarrow$ c) trivial, car $\mathbb K(\lambda)$ est égal à $\mathbb K$ ou est une ext.
quadratique purement imaginaire de $\mathbb K$.

c) $\Longrightarrow$ a) car si $\lambda_\alpha=\sigma_\alpha(\lambda)$, on aura \deleted{\unsure{que} $\bar\lambda_\alpha$ est} \inserted{$\lambda_\alpha\in\mathbb L$, donc}
$\bar\lambda_\alpha$ est l'unique \deleted{\unsure{sa}} \unsure{conjuguée|conjugué} de $\lambda_\alpha$ sur $\mathbb K$, \deleted{\unsure{il}} donc
\[
\lambda_\alpha\bar\lambda_\alpha=N_{\mathbb L/\mathbb K}(\lambda_\alpha)=N_{\mathbb L/\mathbb K}(\lambda)=\lambda\bar\lambda=Q.
\]

\deleted{\uline{Corollaire} \inserted{Supposons $Q\in\mathbb Q^{*+}$.} Soit $L\subset\mathbb L$, \inserted{$L$ fini sur $\mathbb Q$, $\deleted{\unsure{L}}$ non \deleted{\unsure{réel tot}} réel (i.e. non tot. réel)} Alors $L$ est engendré par un
$\lambda$ tel que $\lambda\bar\lambda=Q$. Si $Q\in\mathbb Z$, on peut supposer même
que $\lambda$ est entier.}

\uline{Remarque} Soit $K$ une extension finie tot. réelle de $\mathbb Q$,
et $Q\in\mathbb Q^{\unsure{*+|+*}}$, $Q>1$. Je dis qu'il existe alors une extension quadr.
purement imaginaire $L$ de $K$, \deleted{telle que} et un
entier $i\ge\unsure{0|1}$, \deleted{tels que} tels que $\deleted{\unsure{K}}L$ soit engendré
\inserted{sur $\mathbb Q$} par un élément $\lambda$ tel que $\lambda\bar\lambda=Q^i$ i.e. $|\lambda|^2=Q^i$.
Si $Q\in\mathbb Z$, on peut $\bar{\mathrm m}$ prendre $\lambda$ entier algébrique.

Dém. Soit $a\in K$ tel que $\deleted{\unsure{L|K}}K=\mathbb Q(a)$. \inserted{On peut supposer $a$ entier.} Soit
$i$ assez grand pour que $\deleted{\unsure{a^2}}4Q^i-a^2$ soit tot. $>0$
(\deleted{\unsure{il est possi}} existe, car $Q\in\mathbb Q$, $Q>1$). Soit $\lambda$ solution
de l'équation $\lambda^2-a\lambda+Q^i=0$. Il satisfait à la
condition voulue; si $Q$ est entier, il est entier, à l'étant.
% archive sheet number: 131
\scanpage{133}
Notons que $\lambda'=\lambda^{1/i}$ sera aussi dans $\mathbb L$ (car $\lambda\bar\lambda=Q\,\unsure{\varepsilon|\in|i}$, et
de $=$ pour ses conjugués\ldots), donc on voit aussi qu'il
existe un $\lambda'$ \deleted{\unsure{entier}} $\in\mathbb L$, avec $\deleted{\unsure{\lambda}}\lambda'\bar\lambda'=|\lambda'|^2=Q$, tel que
$\mathbb Q(\lambda')$ soit \deleted{\unsure{pu|p}} \inserted{(purement)} imaginaire et \uline{contienne} $K$; de plus, si
$Q$ est entier, on peut prendre $\lambda'$ entier.
% paper watermark in outline at centre; paper clip at upper left
% archive sheet number: 132
\scanpage{135}
\uline{Groupes de Galois motiviques commutatifs.}

$k$ corps

$M$ catégorie de (iso)-motifs ($1/2$ simples) sur $k$

$G$ lien de $M$, \deleted{supposé commutatif}.

$\unsure{\mathscr G|\mathcal G}$ gerbe des \deleted{\unsure{les}} foncteurs fibres de $M$

$\deleted{\unsure{G|\mathscr G}}$

$\deleted{\unsure{\mathscr G\longrightarrow|\mathscr G_0\longrightarrow}}\uline{M_0\subset M}$, catégorie \deleted{\unsure{des objets|des Tate}} \inserted{engendrée} par $\mathbb Q(1)$ de Tate

$\mathscr G_0,G_0=G_m$ gerbe et lien correspondants

\diagram{\begin{tikzcd}[ampersand replacement=\&]\mathscr G\arrow[r]\&\mathscr G_0\\G\arrow[r,"\varepsilon","{\text{\uline{épimorphisme}}}"']\&G_m\end{tikzcd}}

$\mathscr G_0$ est trivialisé i.e. muni d'une
section, correspondant \deleted{\unsure{au}} \inserted{à un} foncteur-fibre $\varphi$ \deleted{$\mathbb Q(1)\longrightarrow\mathbb Q$} choisi sur $M_0$
[\uline{NB} Comme $H^1(\mathbb Q,G_m)=0$, deux foncteurs
fibres sur $M_0$ sont \inserted{\deleted{\unsure{can}}} isomorphes]

Une trivialisation de $\mathscr G_0$ permet de \uline{restreindre} \unsure{les|le}
lien de $\mathscr G$ de $G$ à $\operatorname{Ker}\varepsilon=\deleted{\unsure{V|U}}U$ : on prend les
foncteurs fibres \uline{munis} d'un isom. de leur
restriction à $M_0$ avec le foncteur fibre fixé
sur $M_0$. On trouve une nouvelle gerbe
% slanted
\margin{\uline{\unsure{réf|vérif} Giraud ?}}
$\mathscr G_{\varphi}$ de lien $U$.

Soit $M'$ la ss-catégorie de $M$ formée des objets
$M$ tels que les $\overset{\unsure{n|-}}{\otimes}(M+\check M)$ soient combinaisons d'un
nb fini \deleted{\unsure{s}} d'objets simples. \inserted{$\mathscr G'$ la gerbe associée, $G'$ le lien.} On en déduit

\diagram{\begin{tikzcd}[ampersand replacement=\&,column sep=small]
{}\&{}\&\mathscr G\arrow[r]\&\mathscr G'\arrow[r,dash]\&1\\
1\arrow[r,dash]\&G^\circ\arrow[r,dash]\&G\arrow[r,"{\text{épim}}"']\&G'\arrow[r,dash]\&1
\end{tikzcd}}
Le noyau de $G\to G'$ est
la ``composante neutre'' $G^\circ$.
D'autre part, l'on sait que
$M'$ est équiv.\textsuperscript{te} à une ss-catégorie \inserted{pleine et stable} des $\mathbb Q$-\unsure{ensembles|modules|vectoriels} tordus sur $k$. La donnée
d'une clôture algébrique $\bar k$ de $k$ fournit un
% stamp: DR (repeated)
% archive sheet number: 134
\scanpage{137}
foncteur fibre\inserted{$\psi=\psi_{\bar k}$} sur $M^{\prime}$, i.e. une section de $\mathscr G^{\prime}$,
d'où une \deleted{\unsure{catégorie|extension} \unsure{résiduelle|restriction}} gerbe
$\mathscr G_{\psi}$ de lien $G^\circ$, qui \deleted{\unsure{donne}} est \deleted{\unsure{liée}}
reliée à la gerbe $\mathscr G$ de lien $G$ via $G^\circ\subset G$.

\uline{N.B} En général, $H^1(\mathbb Q,\Gamma)\ne 1$, donc il y a
ici un choix \uline{très} \uline{particulier} de trivialisation de
\uline{la gerbe $M'$}.

\margin{\uline{\unsure{motifs}} de \uline{niveau $0$} \uline{dans $M$}}% slanted
Considérons enfin la \deleted{gerbe}\inserted{\uline{ss-catégorie \unsure{tensorielle|tannakienne} $M^{(0)}$ de $M$}}
(\unsure{engendrée} par $M_0$ et $M'$ : c'est la catégorie formée des
sommes directes $\displaystyle\sum_{i\in\mathbb Z}V_i(i)$, où les
$V_i$ sont $\in M'$, et les $V_i\in M'$ sont tous
nuls sauf un nbre fini. $\big[$ N.B $\big(\sum V_i(i)\big)\otimes\big(\sum W_j(j)\big)$
\[
\begin{aligned}
 &=\sum_h\left(\sum_{i+j=h} V_i\otimes W_j\right)(h),&
 \left(\sum V_i(i)\right)^\vee&=\sum V_{\unsure{i|-i}}^\vee(-i)\bigg)
\end{aligned}
\]
On trouve un foncteur fibre
\[
 \sum_{i\in\mathbb Z}V_i(i)\longrightarrow\sum V_i(\bar k).\quad,
\]
\margin{\diagram{\begin{tikzcd}[ampersand replacement=\&]
 M_{\unsure{0|00}}\arrow[r,hook,dash]\& M^{(0)}\\
 \& M'\arrow[u,leftrightarrow]
\end{tikzcd}}}
\deleted{\unsure{et|d'où}} $\mathscr G^{(0)}$, $\underset{\text{\unsure{le groupe}}}{G^{(0)}}\simeq G'\times\mathbb G_m$
\diagram{\begin{tikzcd}[ampersand replacement=\&,column sep=small]
 {}\&{}\&\mathscr G\arrow[r]\&\mathscr G^{(0)}\&{}\\
 1\arrow[r]\&U\arrow[r]\&G\arrow[r]\&G^{(0)}=\mathbb G_m\times G'\arrow[r,dash]\&1
\end{tikzcd}}
On trouve une \deleted{\unsure{catégorie|neutralisation}}\inserted{gerbe} $\mathscr G_{\chi}$ sous $U$,
reliée à la grande $\mathscr G$ sous $G$ \deleted{\unsure{par}} et $U\hookrightarrow G$.
% archive sheet number: 136
\scanpage{138}
\mspage{- 16 -}\hfill n° 505

\uline{P. 129}, ligne 3 du haut, remplacer ``métrisable'' par ``séparé''. Ligne 19 du bas,
remplacer ``\uline{métrisable}'' par ``\uline{séparé}''.

Lignes 2 à 10 du bas, remplacer par :

Considérons en effet l'application surjective $f:\{0,1\}^{\mathbb N}\to I$ qui, à tout
point $(\epsilon_n)_{n\geq 0}$ avec $\epsilon_n=0$ ou $\epsilon_n=1$, fait correspondre le nombre
$\displaystyle\sum_{n=0}^{\infty}\epsilon_n\,2^{-n}$ ayant $(\epsilon_n)$ pour développement dyadique. Soit $D$ l'ensemble dénombrable dans $\{0,1\}^{\mathbb N}$ formé des suites autres que la suite $0$, n'ayant qu'un nombre fini de termes $\ne 0$ ; $P=\{0,1\}^{\mathbb N}-D$ est un espace polonais éparpillé (n° 1,
th. 1), et la restriction de $f$ à $P$ est une \uline{bijection} continue de $P\unsure{'|{}}$ sur $I$
(chap. IV, § 8, n° 6).

\uline{P. 130}, la déf. 7 devient déf. 8 ; après la ligne 4 du haut, insérer :

COROLLAIRE.- \uline{Toute somme dénombrable d'espaces lusiniens est un espace lusinien}.
\uline{Toute limite projective dénombrable} (chap. I, § 4, n° 4) \uline{d'espaces lusiniens est un espace lusinien}.

La première assertion découle aussitôt de la déf. 5 et du n° 1, prop. 1 ; la seconde résulte de ce qu'une limite projective dénombrable d'espaces séparés est un sous-espace fermé du produit (dénombrable) de ces espaces, et de ce qu'on a vu au cours de la démonstration de la prop. 11.

\textsuperscript{*}On peut montrer que les espaces de distributions $D'(\Omega)$, $E'(\Omega)$, etc. sur un ouvert $\Omega$ de $\mathbb R^n$, munis de la topologie forte, sont des espaces lusiniens.\textsubscript{*}

Lignes 16 à 20 du haut, remplacer par : ensemble de la topologie \uline{limite projective} des topologies discrètes sur les $C_n$ (chap. I, § 4, n° 4).

Ligne 10 du bas, remplacer ``\uline{fermées}'' par ``\uline{ouvertes}''. Ligne 6 du bas, ajouter à la fin : et contienne les adhérences $\overline{\varphi_{n+1}(c')}$ des $\varphi_{n+1}(c')$ dans $E$ (pour
% archive sheet number: 137
